\documentclass[10pt,a4paper]{amsart}
\usepackage[margin=19mm]{geometry}
\usepackage{amsmath,amssymb,mathtools}
\usepackage[scr=rsfs,cal=euler]{mathalfa}
\usepackage{xfrac}
\usepackage[unicode,hidelinks,bookmarks=false]{hyperref}
\newcommand{\ie}{i.e.\ }
\newcommand{\cf}{cf.\ }
\newcommand{\resp}{respectively\ }
\newcommand{\Subsection}{\subsection}
\providecommand{\nobreakdash}{\nobreak\hbox{-}}
\setlength{\emergencystretch}{2em}
\multlinegap0pt
\usepackage{amssymb}
\usepackage{mathtools}
\usepackage{xcolor}
\usepackage[shortlabels]{enumitem}
\usepackage{tikz}
\usepackage{float}
\newtheorem{theorem}{Theorem}
\newcommand{\dd}{\,\mathrm d}
\newcommand{\erf}{\operatorname{erf}}
\title{Zonoids whose polars are zonoids:\\ the Banach--Mazur distance need not tend to one}
\author{Dmitry Ryabogin, Artem Zvavitch}
\date{Source: arXiv:2609.10852v1 (CC BY 4.0)}
\begin{document}
\maketitle

\noindent\emph{Source and licence.} Zonoids whose polars are zonoids:\\ the Banach--Mazur distance need not tend to one. Version arXiv:2609.10852v1 (CC BY 4.0), distributed by arXiv under the Creative Commons Attribution 4.0 International licence, http://creativecommons.org/licenses/by/4.0/, as recorded in the arXiv licence metadata for this exact version and shown on the abstract page https://arxiv.org/abs/2609.10852v1. Full licence notice: \texttt{licenses/arxiv-2609.10852-CC-BY-4.0.txt}.\par\smallskip

\noindent\textit{Editorial scope.} Counted unit: the article's theorem thm:complete-monotonicity (source0.tex lines 1277--1283). The article's complete original proof of it is delivered with it (source0.tex lines 1404--1586, one \texttt{proof} environment of 183 lines, which the article places away from the statement). No mathematics is written, repaired or re-derived: the statement and the proof are the article's own lines, in the article's own notation, and no retained line is substituted or re-flowed.  Retained uncounted as the source-local context the proof needs: lines 455--455; lines 1266--1270; lines 1341--1344; lines 1362--1364.  Nothing was omitted from any retained span, so the package carries no omission record.  No retained line contains a citation, so the wrapper carries no bibliography and no bibliography entry.  No retained line contains a figure environment, an image-inclusion command or drawing code, so no image is delivered and the package declares no asset dependency.  Editorial formatting only, in addition: an AMS \texttt{amsart} preamble that restates the article's own declarations and macros used by the retained lines, namely the environments \texttt{theorem} on separate counters, together with the standard packages those lines need.  The retained environments print in this document as Theorem 1 in order of appearance, whereas the article prints the same items with the numbering of its own full document.  The complete native source of the article as distributed in the arXiv e-print for this exact version is bundled unchanged as \texttt{sources/209955/source0.tex}.
\subsection{Complete monotonicity and Bernstein functions}\label{compmonotone}

\begin{equation}\label{eq:G-g-definition}
 G(a):=\sqrt\pi e^{a^2}\erf(a),
 \qquad g(y):=e^{(G^{-1}(y))^2},
 \qquad a,y\geq0.
\end{equation}

\begin{equation}\label{eq:P-definition}
 P(u):=\frac1{1+\sqrt u\,G^{-1}(\sqrt u)},
 \qquad u\geq0.
\end{equation}

\begin{equation}\label{eq:G-inverse-ODE}
 (G^{-1})'(y)=\frac1{2(1+yG^{-1}(y))}.
\end{equation}

\begin{theorem}\label{thm:complete-monotonicity}
For $G$ and $g$ defined by \eqref{eq:G-g-definition}, the function
\begin{equation}\label{eq:Q-definition}
 Q(u):=g''(\sqrt u),\qquad u\geq0,
\end{equation}
is completely monotone.
\end{theorem}

\begin{proof}[Proof of Theorem~\ref{thm:complete-monotonicity}]
We begin by computing $Q(u)=g''(\sqrt{u})$.  Since
\[
 g(y)=e^{(G^{-1}(y))^2},
\]
equation \eqref{eq:G-inverse-ODE} gives
\[
\begin{aligned}
 g'(y)
 &=2(G^{-1})'(y)G^{-1}(y)e^{(G^{-1}(y))^2} \\
 &=\frac{G^{-1}(y)e^{(G^{-1}(y))^2}}
         {1+yG^{-1}(y)}.
\end{aligned}
\]
Differentiating once more and again using
\eqref{eq:G-inverse-ODE}, we obtain
\begin{align*}
 g''(y)
 &=
 \frac{
 e^{(G^{-1}(y))^2}
 \bigl(1+2(G^{-1}(y))^2\bigr)(G^{-1})'(y)}
 {1+yG^{-1}(y)}-
 \frac{
 G^{-1}(y)e^{(G^{-1}(y))^2}
 \bigl(G^{-1}(y)+y(G^{-1})'(y)\bigr)}
 {\bigl(1+yG^{-1}(y)\bigr)^2}
 \\
 &=
 \frac{e^{(G^{-1}(y))^2}}
      {\bigl(1+yG^{-1}(y)\bigr)^2}
 \left[
 (G^{-1})'(y)
 \Bigl(
  1+2(G^{-1}(y))^2+2y(G^{-1}(y))^3
 \Bigr)
 -(G^{-1}(y))^2
 \right]
 \\
 &=
 \frac{e^{(G^{-1}(y))^2}}
      {\bigl(1+yG^{-1}(y)\bigr)^2}
 \left[
 \frac{
  1+2(G^{-1}(y))^2+2y(G^{-1}(y))^3}
 {2\bigl(1+yG^{-1}(y)\bigr)}
 -(G^{-1}(y))^2
 \right]
 \\
 &=
 \frac{e^{(G^{-1}(y))^2}}
      {2\bigl(1+yG^{-1}(y)\bigr)^3}.
\end{align*}
By \eqref{eq:P-definition}, for $y\geq0$,
\[
 \frac1{1+yG^{-1}(y)}=P(y^2).
\]
Consequently,
\begin{equation}\label{eq:Q-P-formula}
 Q(u)=g''(\sqrt{u})
 =\frac12e^{(G^{-1}(\sqrt{u}))^2}P(u)^3.
\end{equation}

Although $P^3$ is completely monotone,
\eqref{eq:Q-P-formula} does not immediately imply that $Q$ is
completely monotone, because the exponential factor is increasing.
Since $P(0)=1$ and $G^{-1}(0)=0$, formula
\eqref{eq:Q-P-formula} gives $Q(0)=1/2$.  Moreover, $Q(u)>0$.
Define
\begin{equation}\label{eq:B}
 B(u):=-\log\frac{Q(u)}{Q(0)}.
\end{equation}
Then $B(0)=0$ and
\[
 B'(u)=-\frac{Q'(u)}{Q(u)}.
\]
It is therefore enough to prove that $B'$ is completely monotone.
Indeed, this would imply that $B'\geq0$ and hence that $B\geq0$, so
$B$ would be a Bernstein function.  Property~(iii) in
Subsection~\ref{compmonotone} would then give
\[
 Q(u)=Q(0)e^{-B(u)}=\frac12e^{-B(u)},
\]
which is completely monotone.  We now express $B'$ in terms of $P$.  By
\eqref{eq:G-inverse-ODE} and \eqref{eq:P-definition}, for $u>0$,
\begin{equation}\label{eq:P-inverse-identities}
 (G^{-1})'(\sqrt{u})=\frac12P(u),
 \qquad
 G^{-1}(\sqrt{u})
 =\frac{1-P(u)}{\sqrt{u}\,P(u)}.
\end{equation}
Differentiating \eqref{eq:P-definition} directly and using
\eqref{eq:P-inverse-identities}, we obtain
\begin{align}
 P'(u)
 &=-P(u)^2\left(
 \frac{G^{-1}(\sqrt{u})}{2\sqrt{u}}
 +\frac12(G^{-1})'(\sqrt{u})
 \right) \notag\\
 &=-P(u)^2\left(
 \frac{1-P(u)}{2uP(u)}+\frac14P(u)
 \right) \notag\\
 &=-\frac{P(u)(1-P(u))}{2u}-\frac14P(u)^3.
 \label{eq:P-derivative}
\end{align}
The same identities give
\begin{equation}\label{eq:G-inverse-square-derivative}
 \frac{\dd}{\dd u}\bigl(G^{-1}(\sqrt{u})\bigr)^2
 =
 \frac{G^{-1}(\sqrt{u})}{\sqrt{u}}\,
 (G^{-1})'(\sqrt{u})
 =\frac{1-P(u)}{2u}.
\end{equation}

We now differentiate \eqref{eq:Q-P-formula} directly.  Using
\eqref{eq:P-derivative} and
\eqref{eq:G-inverse-square-derivative}, we find
\begin{align*}
 Q'(u)
 &=
 \frac12e^{(G^{-1}(\sqrt{u}))^2}P(u)^3
 \frac{1-P(u)}{2u}
 +\frac32e^{(G^{-1}(\sqrt{u}))^2}P(u)^2P'(u)
 \\
 &=
 \frac12e^{(G^{-1}(\sqrt{u}))^2}P(u)^3
 \frac{1-P(u)}{2u}
 \\
 &\qquad
 +\frac32e^{(G^{-1}(\sqrt{u}))^2}P(u)^2
 \left(
 -\frac{P(u)(1-P(u))}{2u}-\frac14P(u)^3
 \right)
 \\
 &=
 -\frac12e^{(G^{-1}(\sqrt{u}))^2}P(u)^3
 \left(
 \frac{1-P(u)}u+\frac34P(u)^2
 \right)
 \\
 &=
 -Q(u)\left(
 \frac{1-P(u)}u+\frac34P(u)^2
 \right).
\end{align*}
Since $Q(u)>0$, it follows that
\begin{equation}\label{eq:Q-log-derivative}
 B'(u)=-\frac{Q'(u)}{Q(u)}
 =\frac{1-P(u)}u+\frac34P(u)^2.
\end{equation}

The product $P^2$ is completely monotone by property~(i) in
Subsection~\ref{compmonotone}.  To treat the other term, the
Hausdorff--Bernstein--Widder theorem gives a positive measure $\rho$
on $[0,\infty)$ such that
\[
 P(u)=\int_{[0,\infty)}e^{-ur}\,\dd\rho(r).
\]
Since
$
 \rho([0,\infty))=P(0)=1,
$
for $u>0$ we have
$$
 \frac{1-P(u)}u
 =
 \int_{[0,\infty)}\frac{1-e^{-ur}}u\,\dd\rho(r) =
 \int_{[0,\infty)}\int_0^r e^{-us}\,\dd s\,\dd\rho(r)=
 \int_0^\infty e^{-us}\rho([s,\infty))\,\dd s,
$$
where the last equality follows from Tonelli's theorem.  Thus
$(1-P(u))/u$ is the Laplace transform of a positive measure and is
therefore completely monotone.

Thus $B'(u)$ is completely monotone.  Moreover, $B(0)=0$ and $B'(u)\geq0$, so
$B(u)\geq0$.  Hence $B$ is a Bernstein function.  Finally,
\eqref{eq:B} gives
\[
 Q(u)=Q(0)e^{-B(u)}=\frac12e^{-B(u)},
\]
which is completely monotone by property~(iii) in
Subsection~\ref{compmonotone}.  This completes the proof.
\end{proof}

\end{document}
